\section{Ride-Sharing Systems}

One of the most revolutionary developments in urban mobility is ride-sharing, which is changing the way people commute by allowing travelers on comparable routes to share a vehicle. Mitropoulos et al. (2021) claim that ride-sharing services act as a middleman between drivers and passengers, providing economical and ecologically friendly travel options while resolving issues with pollution and urban congestion\cite{mitropoulos2021}.

Alonso-González et al. (2020) looked at the elements that influence users' propensity to share rides and discovered that social comfort, cost savings, journey time, and privacy all have a significant impact on participation decisions\cite{alonso2020}. In a similar vein, Gangadharaiah et al. (2023) introduced the Pooled Rideshare Acceptance Model (PRAM), highlighting the importance of convenience, perceived safety, and trust for the long-term uptake of ridesharing services\cite{gangadharaiah2023}.

Even though ride-sharing services like Uber, Lyft, and BlaBlaCar have their benefits, they frequently work under centralised commercial frameworks that put making money first through commissions and spike pricing.Ashkrof et al. (2022) noted that driver acceptance and behavior toward ride offers are also influenced by dynamic pricing and operational flexibility\cite{ashkrof2022}.


\subsection{User Factors and Barriers in Ride-Sharing Adoption}
A major body of research highlights user psychology and behavioral factors as critical determinants in ride-sharing adoption. Merat et al. (2017) examined the human factors and user acceptance challenges in ride sharing, particularly for automated and digitalized vehicles, noting that perceived control and safety significantly influence user willingness to participate \cite{merat2017}.

Abouelela et al. (2022) found that the frequency of pooled versus solo ride-hailing use depends on demographic characteristics, trip purpose, and time flexibility\cite{abouelela2022}. Dastani et al. (2024) extended this by proposing a mathematical model that integrates transfer convenience, waiting time, and social comfort into ride-sharing preference predictions\cite{dastani2024}.

Participation in ride-sharing is also influenced by social and geographic differences. Users in low-income or less digitally connected neighbourhoods are less likely to use pooled ride-hailing, according to Soria and Stathopoulos (2021), underscoring the significance of equitable design in mobility platforms\cite{soria2021}. According to Compostella et al. (2025), future driverless and automated ride-sharing scenarios will increase consumer perceptions of safety, control, and comfort\cite{compostella2025}.


\subsection{Passenger Safety }
Adoption of ride-sharing is still heavily influenced by safety concerns. According to Chaudhry et al. (2018), ride-sharing systems should have emergency warning capabilities, user rating systems, and identity verification to guarantee accountability since passenger safety is a major factor in determining trust and retention\cite{chaudhry2018}.
Similarly, Acheampong et al. (2021) emphasized that digital mobility platforms should incorporate transparent user data handling and behavioral monitoring to mitigate risks related to harassment, fraud, and privacy breaches\cite{acheampong2021}. 

According to Mageto et al. (2025), users' propensity to utilise ride-hailing services is greatly reduced by perceived hazards, such as a lack of driver screening and ambiguity about route safety. The design of Wheel Share, which has an emergency feature, user authentication system, and rating and review systems to promote community trust and safety, is directly influenced by these insights
\cite{mageto2025}.

\subsection{Trust Mechanisms in Ride-Sharing}
Hussain et al. (2025) explored commuters’ preferences for future ride-sharing and concluded that a combination of safety assurance, technological reliability, and social responsibility drives long-term engagement\cite{hussain2025}.

\subsection{Technological Advancements in Ride-Sharing Systems}
In terms of scalability, accessibility, and real-time response, emerging technologies have completely changed how ride-sharing systems are implemented. According to Taiebat et al. (2022), data-driven insights can optimize matching and pricing algorithms by analyzing user behavior and predicting ride-sharing preferences using machine learning models\cite{taiebat2022}.

In order to promote justice while maintaining platform operations, Bujak and Kucharski (2024) emphasised the significance of striking a balance between profit and traveler acceptance in ride-pooling through customised fare schemes. This idea can be used to community-oriented systems like Wheel Share
\cite{bujak2024}.

\subsection{Automated Ride Monitoring & Safety Insights}
In 2023, R. Taylor et al. presented the Experian Hunter system, originally designed to analyze large datasets and identify unusual patterns. In a ride-sharing context, similar analytical techniques can help understand user behavior,improve platform reliability, and support better decision making while ensuring data privacy and security\cite{taylor2023}.


\subsection{Experian Hunter}
The Experian Hunter system, a data-driven analytical tool initially created to evaluate and analyse massive data sets for anomalous patterns, was introduced by R. Taylor et al. in 2023. Similar analytical techniques can be applied to a ride-sharing setting to examine ride patterns, user activity, and platform interactions in order to guarantee consistency, boost operational effectiveness, and facilitate trustworthy decision-making\cite{taylor2023}. These tools assist ride-sharing services improve service quality and uphold operational and data protection regulations by offering immediate insights and predictive capabilities. 


\subsection{Intelligent Ride Pattern Monitoring}
In 2025, J. Smith et al. discussed the Actimize system, which uses artificial intelligence and predictive analytics to identify unusual patterns and optimize operational decision making\cite{smith2025}. Similar AI-based analytical techniques can be adapted for ride-sharing platforms to improve user safety, ensure reliable matching, and strengthen overall platform trust.
By analyzing behavioral patterns including unusual ride cancellations, duplicate ride postings, or repeated payment disputes predictive models can identify potential misuse or suspicious activity. Real-time alerts and automated risk scoring can then notify administrators for early intervention.

\subsection{FICO Falcon Fraud Manager}
n 2022 A.Lee et al enhanced that FICO Falcon Fraud Manager operates through machine learning to analyze insurance payments for detecting irregular patterns through its artificial intelligence capabilities. Insured party conduct data alongside existing claim statistics combined with payment documents allows FICO Falcon Fraud Manager to discover potential fraudulent activities. Insurers benefit from real-time fraud detection through the system which develops security systems and minimizes false claims by offering data-based choices to reduce insured losses\cite{lee2022}.


\subsection{ClearSale}
In 2019 B. Johson et al presented that ClearSale combines artificial intelligence with human expertise to use its fraud prevention platform for insurance claims verification so it can stop fraudulent activities. Machine learning algorithms analysis within the platform examines claims data to discover system abnormalities that result in assessment report generation. The user-invoked review process improves both fraud detection quality and enables rapid and seamless handling of genuine claims\cite{johnson2019}.

\subsection{Comparison Table}

\renewcommand{\arraystretch}{1.3}
\setlength{\tabcolsep}{4pt}
\small

\begin{longtable}{|c|c|l|
>{\raggedright\arraybackslash}p{3.2cm}|
>{\raggedright\arraybackslash}p{3.5cm}|
>{\raggedright\arraybackslash}p{3.5cm}|}


\hline
\textbf{Ref} & \textbf{Year} & \textbf{Author(s)} &
\textbf{Approach} & \textbf{Contribution} & \textbf{Limitation} \\
\hline
\endfirsthead

\hline
\textbf{Ref} & \textbf{Year} & \textbf{Author(s)} &
\textbf{Approach} & \textbf{Contribution} & \textbf{Limitation} \\
\hline
\endhead

\hline
\multicolumn{6}{r}{\textit{Continued on next page}} \\
\hline
\endfoot

\hline
\endlastfoot

\cite{mitropoulos2021} & 2021 & Mitropoulos et al. &
Study of ride-sharing’s role in urban mobility &
Highlights cost effectiveness and environmental benefits of shared commuting &
Limited focus on safety and trust mechanisms \\ \hline

\cite{alonso2020} & 2020 & Alonso-González et al. &
Behavioral analysis of ride-sharing adoption &
Identifies cost, time, and privacy as key adoption factors &
Does not address fraud or security issues \\ \hline

\cite{gangadharaiah2023} & 2023 & Gangadharaiah et al. &
Pooled Rideshare Acceptance Model (PRAM) &
Establishes safety, trust, and convenience as central to adoption &
Focused mainly on user perception, not implementation \\ \hline

\cite{chaudhry2018} & 2018 & Chaudhry et al. &
Passenger safety in digital mobility &
Proposes safety and emergency alert mechanisms for trust building &
Does not consider real-time fraud or misuse detection \\ \hline

\cite{acheampong2021} & 2021 & Acheampong et al. &
Data transparency and accountability in ride-sharing &
Promotes secure data handling and user behavior monitoring &
Focus limited to privacy rather than fraud prevention \\ \hline

\cite{taylor2023} & 2023 & R. Taylor et al. &
Experian Hunter automated Ride Monitoring &
Safety insights analytics for real-time risk assessment &
Initially designed for safety; needs adaptation for ride-sharing data \\ \hline

\cite{smith2025} & 2025 & J. Smith et al. &
Intelligent ride pattern monitoring and predictive analytics &
Detects anomalies through behavioral and transactional data patterns &
Industry-specific model; lacks mobility-related parameters \\ \hline

\cite{lee2022} & 2022 & A. Lee et al. &
FICO Falcon Fraud Manager using machine learning &
Provides real-time analysis of payment irregularities &
Limited to financial datasets; not optimized for ride-sharing context \\ \hline

\cite{johnson2019} & 2019 & B. Johnson et al. &
ClearSale AI and human review hybrid model &
Enhances fraud detection accuracy through human-in-the-loop verification &
Dependent on manual review, which may reduce scalability \\ \hline

\cite{hussain2025} & 2025 & Hussain et al. &
Trust mechanisms in ride-sharing systems &
Identifies safety and social responsibility as long-term engagement drivers &
Conceptual focus; lacks technical model implementation \\ \hline
\caption{Summary of Studies on Ride-Sharing}
\end{longtable}
\normalsize
\clearpage



